% Generated draft: requires mathematical and source review.
\documentclass[11pt]{article}
\usepackage[margin=25mm]{geometry}
\usepackage{fontspec}
\setmainfont{Latin Modern Roman}
\usepackage{amsmath,amssymb,mathtools,unicode-math}
\setmathfont{Latin Modern Math}
\usepackage{enumitem,needspace,xurl,hyperref}
\setlength{\parindent}{0pt}
\setlength{\parskip}{5pt}
\setlength{\emergencystretch}{3em}
\newcommand{\N}{\mathbb{N}}
\newcommand{\Z}{\mathbb{Z}}
\newcommand{\Q}{\mathbb{Q}}
\newcommand{\R}{\mathbb{R}}
\newcommand{\C}{\mathbb{C}}
\newcommand{\F}{\mathbb{F}}
\newcommand{\A}{\mathbb{A}}
\newcommand{\PP}{\mathbb{P}}
\newcommand{\E}{\mathbb{E}}
\newcommand{\eps}{\varepsilon}
\newcommand{\dd}{\,\mathrm{d}}
\begin{document}
\section*{UnsolvedMath ID: 801047}
\textbf{Source problem code:} catalogue entry
\subsection{Definitions and mathematical statement}
The source concerns numerical methods for the Painlev\'e equations, regarded there as integrable ordinary differential equations. Here an isomonodromic formulation means an associated linear differential system whose generalized monodromy data remain fixed under deformation; its associated Riemann--Hilbert formulation characterizes a solution through an analytic boundary-value problem. An integrable discretization means a discrete equation retaining an appropriate integrable structure and approaching the associated differential equation in a suitable continuum limit. The source does not specify the integrable structure or the sense of this limit.

\textbf{Research problem (Clarkson, 2019, Problem 1).} Investigate the following three related questions.
\begin{enumerate}
\item Can numerical methods exploiting integrability improve on standard Runge--Kutta methods and their variants for solving Painlev\'e equations?
\item How can the isomonodromic method, through an associated Riemann--Hilbert problem, be used to develop software for numerical study of Painlev\'e equations?
\item Should integrable discretizations be used for numerical study of Painlev\'e equations, given the existence, asserted in the source, of discrete Painlev\'e equations tending to associated continuous Painlev\'e equations in appropriate continuum limits?
\end{enumerate}

These are methodological research questions rather than a specified theorem or quantified numerical bound. The supplied record does not select a Painlev\'e equation, parameter values, real or complex independent-variable domain, initial or boundary data, or treatment of singularities. It also does not define what constitutes an improvement over Runge--Kutta methods, prescribe a Riemann--Hilbert formulation, or select a discrete equation and continuum scaling. Consequently, no universal performance claim or particular convergence assertion is specified by this draft; a precise numerical task requires these choices to be stated explicitly.
\subsection{Short English statement}
Clarkson asks whether integrability can improve numerical computation of Painlev\'e equations, how associated Riemann--Hilbert problems can inform numerical software, and whether integrable discrete Painlev\'e equations should serve as discretizations. The source leaves the equations, data, domains, comparison criteria, and continuum limits unspecified.
\subsection{Sources}
\begin{itemize}
\item Local UnsolvedMath registry, AMR Open Problem Lists: permanent ID \(801047\), exact source problem code \texttt{catalogue entry}, title \emph{Open Problems for Painleve Equations --- Problem 1}. The supplied record attributes the problem to Clarkson, with proposed year 2019, and labels its status \texttt{open}; this label is supplied curation, not an independently verified literature assessment.
\item Clarkson, \emph{Open Problems for Painlev\'e Equations} (2019), Problem 1. Supplied source URL: \url{https://arxiv.org/abs/1901.10122}. The provenance records access on 2026-10-16 and extraction by \texttt{arxiv\_latex}; these are supplied provenance details, not a claim of independent access or verification.
\end{itemize}
\end{document}
